\section{Discussion}
\label{sec:discussion}

\todo{Write once the backbone sweep and the robustness check are complete. The
argument differs substantially between outcomes; sketches of both are retained
below so that the framing is a deliberate choice rather than a reaction to the
numbers.}

\paragraph{If learned representations do not clear the baseline.}
The claim would be that pretrained molecular embeddings, as commonly used,
carry little information about docking score beyond what twelve
physicochemical descriptors already provide --- and that reported successes
may partly reflect the absence of such a control together with metrics that do
not account for molecular size. The practical implication is immediate: a
surrogate for ultra-large screening need not store terabytes of embeddings,
and the storage freed is better spent docking more molecules. The claim would
require the robustness check of Section~\ref{sec:results-robust} to hold across
model families, and replication on at least one further target.

\paragraph{If they clear it comfortably.}
The contributions become the size-controlled protocol itself, the demonstration
that a substantial share of apparent performance is size-driven, and a
cost--accuracy analysis over embedding dimensionality. The relevant question
shifts to whether the highest-dimensional representations justify their
storage: GROVER occupies \num{6.6} times the bytes per molecule of Uni-Mol v1,
and a comparable accuracy at that ratio is a decision-relevant finding for
anyone provisioning a screening campaign.

\paragraph{In either case.}
Two observations appear robust to the outcome. First, the descriptor baseline
is strong, cheap and non-trivial once size is controlled for, and its absence
from the surrogate literature is a gap worth closing regardless of which side
of it learned representations fall. Second, unweighted late fusion is unsafe
when member quality is uneven; conclusions drawn from benchmarks of comparable
modalities do not transfer automatically to settings where some
representations are much weaker than others.
